\PassOptionsToPackage{table}{xcolor}
% ===================== CONFIGURACIÓN DE PÁGINA MEJORADA =====================
\ifdefined\pdfminorversion\pdfminorversion=7\fi
\usepackage[top=4.5cm, bottom=2.5cm, left=2.0cm, right=2.0cm,
headheight=4cm, headsep=0.5cm]{geometry}

% ===================== TIPOGRAFÍA Y CODIFICACIÓN =====================
\usepackage[utf8]{inputenc}
\usepackage[T1]{fontenc}
\usepackage[spanish,es-tabla]{babel}
\usepackage{tgtermes}
\usepackage{microtype}

% Equivalencias para los símbolos Unicode presentes en la plantilla al usar XeTeX.
\ifdefined\XeTeXversion
\usepackage{newunicodechar}
\newunicodechar{°}{\textdegree}
\newunicodechar{²}{\textsuperscript{2}}
\newunicodechar{³}{\textsuperscript{3}}
\newunicodechar{·}{\ensuremath{\cdot}}
\newunicodechar{±}{\ensuremath{\pm}}
\fi


% ===================== ELIMINAR SANGRÍAS =====================
\setlength{\parindent}{0pt}
\setlength{\parskip}{6pt}

% ===================== CONFIGURACIÓN DE TÍTULOS =====================
\usepackage{titlesec}
\titleformat{\section}{\bfseries\fontsize{12}{14}\selectfont}{\thesection}{1em}{}
\titleformat{\subsection}{\bfseries\fontsize{12}{14}\selectfont}{\thesubsection}{1em}{}
\titleformat{\subsubsection}{\bfseries\fontsize{12}{14}\selectfont}{\thesubsubsection}{1em}{}

% ===================== PAQUETES MATEMÁTICOS =====================
\usepackage{amsmath,amssymb,amsfonts,mathtools}
\usepackage{textgreek,gensymb,siunitx}

% ===================== QUÍMICA Y SÍMBOLOS =====================
\usepackage[version=4]{mhchem}
\usepackage{chemfig,textcomp}

% ===================== TABLAS Y FIGURAS =====================
\usepackage{booktabs,array,longtable,multirow}
\usepackage{graphicx,float,subcaption}
\usepackage{tabularx}

% ===================== SOPORTE PARA SVG =====================
\usepackage{svg}

% ===================== GRÁFICAS VECTORIALES (PGFPLOTS) =====================
% Pipeline: Python exporta CSV a assets/ → \addplot table en las secciones.
\usepackage{pgfplots}
\pgfplotsset{compat=1.18}
\usepackage{pgfplotstable}

% ===================== ENLACES E HIPERVÍNCULOS =====================
\usepackage[hidelinks,colorlinks=true,linkcolor=blue,citecolor=blue]{hyperref}
\usepackage{url}

% ===================== OTROS PAQUETES =====================
\usepackage{enumitem,xcolor,listings,algorithm,algpseudocode}
\usepackage{tocloft}
\usepackage{appendix}
\usepackage{eso-pic}
\usepackage{lastpage}

% ===================== PAQUETES ADICIONALES PARA TABLAS ESTÉTICAS =====================
\usepackage[table]{xcolor} % Necesario para rowcolor si no está en documentclass
\usepackage{array}
\usepackage{colortbl}
\usepackage{multirow}
\usepackage{booktabs}

% ===================== AUTO-NUMERACIÓN DE TABLAS =====================
\newcounter{itemcount}
\newcolumntype{N}{>{\stepcounter{itemcount}\theitemcount}c}

% ===================== CONFIGURACIÓN DEL ENCABEZADO =====================
\usepackage{fancyhdr}

% ===================== INFORMACIÓN DEL DOCUMENTO =====================
% ===================== DATOS DEL PROYECTO =====================
% Toda la información del proyecto se centraliza aquí para facilitar
% la creación de nuevos informes sin alterar el código fuente.

% 1. Información General del Proyecto
\newcommand{\projecttitle}{DIAGNOSTICO DEL SISTEMA DE TRATAMIENTO Y DISTRIBUCIÓN DE AGUA PURIFICADA Y DISTRIBUCIÓN JGB}
\newcommand{\documentcode}{W2608PRINF001}
\newcommand{\documenttype}{INFORME TÉCNICO}
\newcommand{\targetcompany}{JGB}
\newcommand{\projectcity}{Cali}

% 2. Información de Documento y Emisión
\newcommand{\docRevision}{0}
\newcommand{\docFecha}{\today}
\newcommand{\docFechaLarga}{19/10/2026}

% 3. Integrantes y Control de Firmas
\newcommand{\directorProyecto}{Javier Micolta}
\newcommand{\ingenieroPrincipal}{J.Micolta}
\newcommand{\ingenieroEspecialista}{MSc Jonathan Arboleda}
\newcommand{\firmaElaboro}{J.Arboleda}
\newcommand{\firmaReviso}{J.Micolta}
\newcommand{\firmaAprobo}{J.Micolta}

% 4. Fechas de Firmas y Revisiones (formato dd/mm/aaaa según guía de codificación del proyecto)
\newcommand{\fechaElaboro}{dd/mm/aaaa}
\newcommand{\fechaReviso}{dd/mm/aaaa}
\newcommand{\fechaAprobo}{dd/mm/aaaa}

% 5. Datos de Historial de Revisiones (Portada / Firmas)
% Se recomienda llenar según la configuración del proyecto (dd/mm/aaaa)
\newcommand{\fechaRevCero}{dd/mm/aaaa}
\newcommand{\descRevCero}{Borrador de plantilla}
\newcommand{\fechaRevUno}{--}
\newcommand{\descRevUno}{Revisión cliente}
\newcommand{\fechaRevDos}{--}
\newcommand{\descRevDos}{Emisión final}
\newcommand{\fechaRevTres}{--}
\newcommand{\descRevTres}{--}
\newcommand{\fechaRevCuatro}{--}
\newcommand{\descRevCuatro}{--}

% Convención local WXXXX autorizada por el usuario; norma fuente guía de codificación del proyecto conservada.
\newcommand{\projectcode}{W2608}


% Adaptaciones para el encabezado corporativo
\newcommand{\documentrevision}{REV \docRevision}
\newcommand{\documentdate}{\docFechaLarga}
\newcommand{\elaboratedby}{\MakeUppercase{\firmaElaboro}}
\newcommand{\reviewedby}{\MakeUppercase{\firmaReviso}}
\newcommand{\approvedby}{\MakeUppercase{\firmaAprobo}}

% ===================== CONFIGURACIÓN DE LISTINGS =====================
\lstset{
	basicstyle=\footnotesize\ttfamily,
	breaklines=true,
	frame=single,
	language=Python,
	showstringspaces=false
}

\journal{Revista de Ingeniería}